\documentclass[a4paper, 12pt]{extarticle}

\input{../../slides/common/slide-plan.tex}

\usepackage[utf8]{inputenc}
\usepackage[russian]{babel}

\usepackage{enumerate}

\usepackage{hyperref}
\usepackage{multirow} 
\usepackage{graphicx}
\usepackage{bm}
\usepackage{geometry}
\usepackage{listings}
\usepackage{tabularx}
\usepackage{tikz}
\usetikzlibrary{automata,positioning}

\geometry{a4paper,top=1.5cm,bottom=1.5cm,bindingoffset=0cm}
\geometry{left=2cm,textwidth=17cm}
% \linespread{1.0}
  
\usepackage{amssymb}
\usepackage{amsmath}
\usepackage{amsthm}
\usepackage{verbatim}

\title{}
\author{}
\date{}

\begin{document}

\section*{Основные определения, используемые в курсе "Параллельное программирование"}

\section*{Предварительные замечания}

Некоторые концепции и определения переиспользуются из ранее пройденных курсов ''Операционные системы'' и ''Архитектура ЭВМ''.
В таком случае детальное определение помнить наизусть не требуется, просто не путайте потоки и процессы между собой.

Многие определения в курсе ''Параллельное программирование'' вводятся дважды или трижды -- сначала в простом, ''наивном'' варианте, а затем в более зубодробительной формальной форме. Желательно, чтобы вы понимали идею, стоящую за определением (''зачем?'', мотивация), а также умели его применять на практике (''как?'', техническое исполнение).

\section{Блок №1}

\begin{itemize}
	\item \textbf{Конкурентность (concurrency)}. Предположим, многоагентная система состоит из задач и одного исполнителя.
	\begin{itemize}
		\item Любую задачу можно прервать в произвольный момент времени (tasks are interruptible)
		\item Исполнитель может исполнять любую из прерванных задач
		\item В любой фиксированный момент времени исполнитель выполняет не более одной задачи
	\end{itemize}
	В таком случае мы говорим, что система допускает \textbf{конкурентное} исполнение. Благодаря частой смене исполняемых задач, можно добиться иллюзии одновременного исполнения.

	\item \textbf{Параллелизм (parallelism)}. Предположим, многоагентная система состоит из задач и нескольких исполнителей.
	\begin{itemize}
		\item В один и тот же момент времени над разными задачами могут работать разные исполнители (tasks are independent)
		\item Любая задача может исполняться любым из исполнителей
		\item В любой фиксированный момент времени исполнитель выполняет не более одной задачи		
	\end{itemize}
	В таком случае мы говорим, что система допускает \textbf{параллельное} исполнение. Наблюдаемая одновременность достигается с помощью использования нескольких вычислителей.

	\item \textbf{Закон Амдала (Amdahl's law)}. Пусть необходимо решить некоторую вычислительную задачу, состоящую из последовательной части ($\alpha$) и параллельной части ($1-\alpha$). Тогда ускорение, которое может быть получено на вычислительной системе из $p$ процессоров, по сравнению с однопроцессорным решением не будет превышать величины 
	$$S_p=\frac{1}{\alpha+\frac{1-\alpha}{p}}$$

	\item \textbf{Round-robin scheduling}. Алгоритм распределения нескольких задач или/и нагрузки распределённой вычислительной системы. Пусть имеется $N$ процессоров и $M$ задач. Первая задача ($m = 1$) назначается для выполнения первому процессору ($n = 1$), вторая — второму и т. д., до достижения последнего процессора ($m = N$). Тогда следующая задача ($m = N+1$) будет назначена снова первому процессору и т. п. Таким образом происходит перебор выполняющих задания процессоров по циклу, или по кругу (round). Обычно ''задачей'' выступает квант планирования (например, 10 миллисекунд).

	\item \textbf{Переключение контекста (context switch)}. Процедура прекращения выполнения процессором одной задачи (процесса или потока) с сохранением всей необходимой информации и состояния, необходимых для последующего продолжения с прерванного места, и восстановления и загрузки состояния задачи, к выполнению которой переходит процессор.

	\item \textbf{Квант планирования (scheduling quantum)}. Характерный промежуток времени, во время которого, ''в среднем'', процессор не совершает смены контекста. На самом деле это динамически (пере-)вычисляемая величина, тесно связанная с алгоритмом планирования, текущей загрузкой системы и многими другими факторами.

	\item \textbf{Вытесняющая многозадачность (pre-emptive multitasking)}. Вид многозадачности, при которой операционная система принимает решение о переключении между задачами по истечении некоего кванта времени и ''форсирует'' смену контекста, независимо от того, в какой точке исполнения находится пользовательский процесс.

	\item \textbf{Кооперативная многозадачность (cooperative multitasking)}. Вид многозадачности, при котором следующая задача выполняется только после того, как текущая задача явно объявит себя готовой отдать процессорное время другим задачам.

%	\item \textbf{Scheduling metrics}.
%	\begin{itemize}
%		\item utilization
%		\item latency
%		\item efficiency
	%\end{itemize}

	\item \textbf{Процесс (process)}. Абстракция операционной системы. ''Нечто'', представляющее выполняющуюся программу и все её элементы: адресное пространство, глобальные переменные, регистры, стек, открытые файлы и т.д. 

	\item \textbf{Поток (thread)}. Абстракция операционной системы. Наименьшая единица ''работы'', которой оперирует планировщик ОС. Несколько потоков выполнения могут существовать в рамках одного и того же процесса и совместно использовать ресурсы, в первую очередь -- одно и то же виртуальное адресное пространство.
	
	\item \textbf{Threading models}.
	\begin{itemize}
		\item \textbf{1:1}. Каждой ''умозрительной'' сущности соответствует одна ''более практическая'' сущность. Примеры:
		\begin{itemize}
			\item 1 Java поток -- 1 поток операционной системы
			\item 1 поток операционной системы -- 1 ядро процессора
			\item 1 http соединение -- 1 Java поток, который его обрабатывает
		\end{itemize}
		
		\item \textbf{N:1}. Все логические единицы планирования ''мультиплексируются'' на один и тот же ресурс. Примеры:
		\begin{itemize}
			\item все процессы и потоки операционной системы -- 1 ядро процессора (актуально для одноядерных компьютеров)
			\item все легковесные задачи (\texttt{Callable}) вычисляются одним исполнителем (например, \texttt{Executors.newSingleThreadExecutor()})			
		\end{itemize}

		\item \textbf{N:M}. Гибридная схема. Примеры:
		\begin{itemize}
			\item все процессы и потоки операционной системы -- N процессоров.
			\item все легковесные задачи (\texttt{Callable}) вычисляются фиксированным множеством исполнителей (например, \texttt{Executors.newFixedThreadPool()})			
		\end{itemize}

	\end{itemize}

	\item \textbf{Interleaving model}. Модель исполнения многопоточных программ, которая предполагает, что все операции выполняются одним исполнителем (N:1 threading model), исполнение каждого потока может быть прервано в любой точке (с точностью до неделимых операций) и на произвольный промежуток времени. Предполагается, что любое изменение разделяемой памяти сразу становится видно всем потокам исполнения. В данной модели ''время'' измеряется не секундами, а событиями, которые выстраиваются в некоторую последовательность. Такую последовательность мы называем трассой исполнения многопоточной программы (\textbf{concurrent execution trace}). Если программа использует более одного потока, то каждая точка прерывания порождает множество потенциальных трасс исполнения и, как следствие, различных результатов работы программы. 

	\item \textbf{Состояние гонки (race condition)}. Свойство многопоточной программы иметь более одной трассы исполнения. Наблюдаемый недетерминизм поведения и/или результатов может быть как нежелательным дефектом, так и допустимым поведением.

	\item \textbf{Гонка данных (data race)}. Дефект в многопоточных программах. Происходит, если не меньше двух потоков одновременно (без должной синхронизации) обращаются к одному и тому же объекту (ячейке памяти), причем как минимум один из участников меняет состояние объекта (записывает значение). \textbf{Замечание}: если в ячейке памяти хранилось значение \texttt{1} и один из потоков пишет в эту ячейку памяти тоже \texttt{1} -- это \textbf{всё равно} гонка данных.
	
	\item \textbf{Блокирующий метод (blocking method)}. Метод, приостанавливающий исполнение текущего потока до тех пор, пока не наступит некоторое внешнее событие.

	\item \textbf{Wait-for graph}. Направленный граф, в котором вершины соответствуют потокам исполнения программы, а дуги отражают свойство ''поток A ожидает поток B''. По мере исполнения программы и при наступлении каких-то внешних событий wait-for graph динамически изменяется, т.е. это абстракция, которая зависит от текущего состояния вычислений. \textbf{Замечание}: в общем случае, не только блокирующие методы приводят к возникновению дуг; если поток A в цикле ожидает смены флага, который может быть изменен только потоком B -- это тоже считается wait-for дугой.

	\item \textbf{Состояние тупика (deadlock)}. Состояние многопоточной программы, когда её wait-for graph содержит цикл. Иными словами, в программе образовалось некоторое множество потоков, которые ''ждут'' друг друга и потому никогда не будут совершать какой-либо полезной работы.
	
	\item \textbf{Инверсия приоритетов (priority inversion)}. Состояние многопоточной программы, когда скорость исполнения высокоприоритетного потока A зависит от скорости исполнения низкоприоритетного потока B. Например, когда A ждет освобождения мьютекса, который захватил B.

%	\item \textbf{Naive thread-safety}

%	\item \textbf{Visibility}

	\item \textbf{Точка синхронизации (synchronization point)}. Волшебное место в программе, которое запрещает компилятору переставлять операции, процессору кешировать старые значения, а языку программирования вести себя контринтуитивно. В блоке №2 магия будет разбиваться о суровый соленый юмор разработчиков процессоров и дизайнеров языков программирования.

	\item \textbf{Голодание (starvation)}. Свойство многопоточной программы. Некоторый поток, пытающийся воспользоваться многопоточным объектом, может исполнить неограниченное количество инструкций прежде чем получит доступ к данному объекту, при условии, что другие потоки также пытаются воспользоваться этим объектом. Лучше всего вспомнить сказку про Алиса и Боба из второй лекции. 

	\item \textbf{Взаимное исключение (mutual exclusion)}. В любой момент времени, указанный фрагмент программного кода исполняется не более чем одним потоком программы. \textbf{Замечание}: недостижимый код удовлетворяет этому свойству, этого свойства \textbf{недостаточно} для мьютекса.

	\item \textbf{Отсутствие тупиков (deadlock freedom)}. Если несколько потоков пытаются захватить мьютекс, как минимум один из них рано или поздно преуспеет. \textbf{Замечание}: пустая функция удовлетворяет этому свойству, этого свойства \textbf{недостаточно} для мьютекса.

	\item \textbf{Отсутствие голодания (starvation freedom)}. Если поток пытается захватить мьютекс, то рано или поздно он преуспеет.
	\textbf{Замечание}: рано или поздно весьма растяжимое понятие, мьютекс со свойством \textbf{starvation freedom} необязательно практичен.

	\item \textbf{Мьютекс (mutex)}. Примитив синхронизации, запрещающий одновременное исполнение некоторого блока кода. Должен гарантировать mutual exclusion и deadlock freedom. Может (но не обязан) гарантировать starvation freedom, такие разновидностью мьютексов часто называют честными (fair mutex).

	\item \textbf{Критическая секция (critical section)}. Набор инструкций (блок кода), который обрамлен методами \texttt{mutex.lock} и \texttt{mutex.unlock}.

	\item \textbf{Реентерабельность (reentrancy)}. Свойство мьютекса, разрешающее повторный вызов метода \texttt{lock} без самоблокировки потока. Подразумевает, что реализация мьютекса ''считает'' количество входов и выходов в критическую секцию, мьютекс разблокируется (разрешает другим потокам войти в критическую секцию) только когда такой ''логический счетчик'' падает до нуля.

	\item \textbf{Честность (fairness)}. Свойство многопоточного объекта или алгоритма давать всем участникам ''равный'' доступ (например, успешно захватывать мьютекс). Часто сопряжено с повышенной отзывчивостью (responisveness), малым временем отклика (latency), сниженной пропускной способностью (throughput).

	\item \textbf{Admission policy}. Политика ''входа'' в многопоточную структуру данных (мьютекс, монитор, условную переменную).
	Чаще всего описывается в терминах очередности -- ''первый вошел -- первый вышел'' (FIFO, очередь), ''последний пришел -- первый вышел'' (LIFO, стек) и их производных (очередь с приоритетами, рандомизированные алгоритмы).

	\item \textbf{Code locking}. Шаблон многопоточного программирования. Экземпляр мьютекса используется для защиты некоторого блока кода от исполнения несколькими потоками одновременно. В большинстве случаев используется статический мьютекс (единственный экземпляр, singleton).

	\item \textbf{Data locking}. Шаблон многопоточного программирования. Экземпляр мьютекса используется для защиты некоторого набора данных от одновременного доступа, совершаемого различными потоками. В большинстве случаев для независимых данных используются независимые мьютексы.

	\item \textbf{Lock splitting}. Шаблон многопоточного программирования. Экземпляр мьютекса, защищающий некоторый логический модуль программы, заменяется множеством мьютексов. При этом происходит разбиение модуля на части, требуется особая внимательность, чтобы результирующая программа была корректной и не приводила к состоянию тупика (deadlock).

	\item \textbf{Активное ожидание (spinning) с выдержкой (backoff)}. Цикл активного ожидания с проверкой условия (polling) может быть ''разбавлен'' переводом потока в неактивное состояние (например \texttt{Thread.sleep}) 
	\begin{itemize}
		\item ни при каких условиях (no backoff) -- гарантирует максимальное время отклика ценой повышенных расходов CPU
		\item с равной периодичностью (fixed backoff)
		\item c экспоненциально возрастающим периодом ожидания (exponential backoff)
	\end{itemize}

	\item \textbf{Условная переменная (condition variable)}. Примитив синхронизации, позволяющий потоку перейди в неактивное состояние и, при получении сигнала от другого потока, вернуться к исполнению. В большинстве современных многопоточных систем рекомендуется (или форсируется) использование условной переменной в паре с мьютексом.

	\item \textbf{Lost signal}. Дефект в многопоточных программах. Чаще всего происходит из-за логической ошибки: поток A пытается передать некоторую информацию потоку B (например посылает сообщение или вызывает метод \texttt{signal}), а поток B находится в ''невосприимчивом'' состоянии (не проверяет очередь сообщений, не заблокирован в методе \texttt{await}).

	\item \textbf{Spurious wakeup}. Свойство некоторых реализаций блокирующих методов. Наиболее известен как ''беспричинный'' выход из метода \texttt{Condition.await}. Требует оборачивать блокирующий метод в цикл \texttt{while} с перепроверкой условия ожидания.

	\item \textbf{Predicate invalidation}. Дефект в многопоточных программах. Чаще всего происходит из-за логической ошибки: поток выходит из блокирующего метода и ожидает, что некоторый инвариант программы выполнен. При этом в программе возможны сценарии исполнения, когда инвариант для потока не соблюдается.

%	\item \textbf{Thread pool}
%	\begin{itemize}
%		\item single
%		\item fixed
%		\item caching
	%\end{itemize}

	\item \textbf{Resource deadlock}. Дефект в многопоточных программах. Возникает при использовании примитивов синхронизации групп потоков (\texttt{Semaphore}, \texttt{FixedSizeThreadPool}) -- N потоков ожидают, что N+1-ый поток выполнит некоторое ''разблокирующее'' действие, однако этот поток не может получить квант планирования, т.к. все доступные ресурсы (N штук) уже исчерпаны.

	\item \textbf{Monitor}. Примитив синхронизации, комбинирующий в себе мьютекс и единственную условную переменную. Часто имеет языковую поддержку для структурированной синхронизации (\texttt{synchronized} блок в Java, RAII-идиома в C++, \texttt{with} в Python и т.п.).	

	\item \textbf{Wait morphing}. Оптимизация структуры данных \texttt{Monitor} или связанных друг с другом \texttt{Lock+Condition}.
	Нотифицированные потоки не пробуждаются в момент отправления сигнала (это бессмысленно, критическая секция занята отправителем сигнала), а перекладываются в очередь ''готовых к пробуждению''. Только при выходе из критической секции поток-бывший-владелец пробуждает один из готовых потоков.

	\item \textbf{Thundering herd problem}. Дефект в многопоточных программах. Возникает при неуместном использовании механизмов массового пробуждения потоков (broadcast). Множество потоков получает квант планирования, но только часть из них может эффективно им распорядиться (сделать полезную работу), остальные потоки вынуждены остановиться в каком-то блокирующем методе, например при попытке захватить мьютекс.

	\item \textbf{CountDownLatch}. Примитив синхронизации, позволяющий остановить/разблокировать группу потоков.

	\item \textbf{Semaphore}. Примитив синхронизации, позволяющий ограниченному числу потоков находиться внутри некоторого блока кода. 
	
	\item \textbf{ReadWriteLock}. Примитив синхронизации, разбивающий потоки на две группы -- ''читатели'', произвольное число которых может находится в rw-секции одновременно и ''писатели'', среди которых только один может находится в rw-секции. Читатели и писатели никогда не находятся в rw-секции одновременно. \textbf{Замечание}: терминология ''читатели/писатели'' весьма неудачная, потоки первой группы вполне могут записывать какие-то данные, если того требует многопоточный алгоритм. Рекомендуется разобрать алгоритм распределенного счетчика из лекции №5.   
	
	\item \textbf{ThreadLocal}. Потоково-локальные данные. ''Именованная переменная'', значение которой у каждого потока своё, персональное, и потому может изменяться без синхронизации с другими потоками.
	
	\item \textbf{Репликация (replication)}. Шаблон многопоточного программирования. Предполагает использование потоково-локальных переменных для того, чтобы делать копии (реплики) разделяемого состояния. На некоторых путях исполнения реплики используются для эффективных (масштабируемых, неблокирующих) вычислений.
	
	\item \textbf{Снимок состояния (snapshotting)}. Шаблон многопоточного программирования. Предполагает использование неизменяемых (или персистентных) структур данных для того, чтобы минимизировать время исполнения критической секции. 

	\item \textbf{Отмена задач (cancellation)}. Шаблон многопоточного программирования. Отмена еще-не-начатых задач позволяет сэкономить вычислительные ресурсы в системах, которые должны оперативно реагировать на изменяющуюся обстановку.
	\begin{itemize}
		\item \textbf{Forced cancellation}. Отмена уже исполняющихся агентов (потоков, процессов, задач) в середине исполнения может приводить к потере данных и/или невалидному состоянию системы. Не делайте так.
		
		\item \textbf{Cooperative cancellation}. Инициатор отмены посылает сигнал фактическому исполнителю и дожидается ответа (задача была отменена, задача успела завершиться успехом).
	\end{itemize}

	\item \textbf{Cancellation token}. Шаблон многопоточного программирования. Используется для кооперативной отмены задач в многопоточной среде.

	\item \textbf{Concurrent state machine}. Шаблон многопоточного программирования. Предполагает использование потокобезопасного конечного автомата для управления состоянием многопоточной структуры данных. Текущее состояние автомата определяет множество допустимых операций и допустимые переходы в другие состояния.

	\item \textbf{Interruption}. Дизайн многопоточных систем, позволяющий прерывать исполнение потоков в выделенных точках (interruption points). Как правило это ''долгие'' или ''блокирующие'' методы, которые поддерживают безопасный досрочный возврат.

%	\item \textbf{timeout}
%	\item \textbf{concurrent invariant}
%	\item \textbf{assert}
%	\item \textbf{unit testing}
%	\item \textbf{stress testing}
	%\item \textbf{bug probability est}

\end{itemize}

\newpage

\section{Блок №2}

\end{document}
